\section{Níveis de Privilégio}

Os níveis de privilégio surgem com a evolução dos \glspl{so} e com a necessidade de particionar e proteger os recursos de um computador.
Esses níveis são organizados em um modelo de camadas, no qual a camada mais alta possui acesso mais restrito, enquanto a camada mais baixa detém controle total sobre o \gls{hardware}.
Além disso, como os níveis superiores não possuem acesso direto a determinados recursos, os níveis inferiores podem lhes fornecer serviços, abstrações e mecanismos de acesso, sempre de forma controlada.
Embora a implementação desses níveis varie entre diferentes arquiteturas, não é desejável que todas as aplicações executem no nível mais privilegiado, pois isso pode comprometer o isolamento
entre processos e gerar competição descontrolada pelos recursos do sistema.
Por isso, a maioria dos processadores e \glspl{so} adotam, no mínimo, uma divisão binária com o modo usuário e o modo supervisor, ilustrados na Figura~\ref{fig:niveis-privilegio}.

\begin{figure}[htbp]
    \centering
    \begin{tikzpicture}[
        >=Stealth,
        camada/.style={
            draw,
            thick,
            minimum width=6cm,
            minimum height=1.2cm,
            outer sep=0pt
        }
    ]
        \node[camada, fill=blue!5] (u) {Usuário (U-Mode)};
        \node[camada, fill=blue!10, below=0.2cm of u] (s) {Supervisor (S-Mode)};
        \node[camada, fill=blue!15, below=0.2cm of s] (h) {\Gls{hardware}};

    \end{tikzpicture}
    \caption{Exemplo Básico de Níveis de Privilégio.}
    \label{fig:niveis-privilegio}
\end{figure}

O modo usuário é destinado à execução da maioria das aplicações. Nesse nível, os programas operam com acesso restrito aos recursos do sistema, sem executar instruções privilegiadas
nem acessar diretamente componentes críticos da máquina.

Já o modo supervisor é destinado à execução do próprio \gls{so}, dispondo de acesso total aos recursos do sistema. Nesse nível, o sistema operacional gerencia memória, dispositivos
e processos, atuando como intermediário entre os programas em modo usuário e o hardware.


